\section*{Problem 1}

在全连接神经网络中，ReLU 激活函数会将输入空间划分为若干个区域。在每一个固定的激活区域内，网络实际上退化为一个仿射映射。设一个两层全连接神经网络定义为
\[
\bm{z}^{(1)} = W_1 \bm{x} + \bm{b}_1,\quad
\bm{h} = \operatorname{ReLU}(\bm{z}^{(1)}),\quad
\bm{z}^{(2)} = W_2 \bm{h} + \bm{b}_2,
\]
其中
\[
W_1 = \begin{bmatrix}
    1 & -1 \\
    2 & 1 \\
    -1 & 2
\end{bmatrix},\quad
\bm{b}_1 = \begin{pmatrix} 0 \\ -1 \\ 1 \end{pmatrix},\quad
W_2 = \begin{bmatrix} 1 & 0 & 2 \\ -1 & 1 & 0 \end{bmatrix},\quad
\bm{b}_2 = \begin{pmatrix} 1 \\ 0 \end{pmatrix},
\]
对输入
\[
\bm{x} = \begin{pmatrix} 2 \\ 1 \end{pmatrix},
\]
完成以下任务：
\begin{enumerate}
    \item 计算 $\bm{z}^{(1)}$、$\bm{h}$、$\bm{z}^{(2)}$。
    \item 判断该输入对应的 ReLU 激活模式。
    \item 在该激活模式保持不变的区域内，写出网络对应的仿射映射。
\end{enumerate}

\section*{Problem 2}

全连接层的一个主要问题是参数量随输入维度迅速增长。卷积层通过局部连接和权重共享显著减少参数量。设输入是一张大小为 $32 \times 32$ 的 RGB 图像，即输入通道数为 $3$。

现在比较以下两种结构：
\begin{itemize}
    \item 结构 A：将图像展平成向量后，接一个含有 $1024$ 个神经元的全连接层。
    \item 结构 B：使用 $16$ 个 $3 \times 3$ 的卷积核，步长为 $1$，无零填充。
\end{itemize}
假设所有层都含有偏置项。完成以下任务：
\begin{enumerate}
    \item 计算结构 A 的参数量。
    \item 计算结构 B 的参数量。
    \item 计算结构 B 的输出特征图大小。
    \item 若将结构 B 的输出展平后接一个 $10$ 类分类器，计算该分类器的参数量。
\end{enumerate}

\section*{Problem 3}

在多分类任务中，神经网络通常先输出 logits，再通过 softmax 转化为类别概率。设一个三分类模型的 logits 由
\[
\bm{z} = W \bm{x} + \bm{b}
\]
给出，其中
\[
W = \begin{bmatrix}
    1 & 0 \\
    0 & 1 \\
    -1 & 0
\end{bmatrix},\quad
\bm{b} = \begin{pmatrix} 0 \\ 0 \\ 0 \end{pmatrix},\quad
\bm{x} = \begin{pmatrix} 1 \\ 0 \end{pmatrix},
\]
真实标签为第一类，即
\[
\bm{y} = \begin{pmatrix} 1 \\ 0 \\ 0 \end{pmatrix}.
\]
softmax 输出定义为
\[
p_i = \frac{\exp(z_i)}{\sum_{j=1}^{3} \exp(z_j)},
\]
交叉熵损失函数为
\[
L = -\sum_{i=1}^{3} y_i \log p_i.
\]
完成以下任务：
\begin{enumerate}
    \item 计算 $\bm{z}$ 和 $\bm{p}$。
    \item 计算损失 $L$。
    \item 利用 $\nabla_{\bm{z}} L = \bm{p} - \bm{y}$，计算 $\nabla_W L$ 和 $\nabla_{\bm{b}} L$。
\end{enumerate}
